\section{Support and specialization}\label{sec:specialization}

Let $\mathscr X\to B$ be a flat projective morphism from a smooth
complex algebraic fourfold to a smooth curve, with central fiber $X_0$
having only ordinary double points. Suppose a nearby fiber $S$ and a
projective small resolution $p:Y\to X_0$ are smooth Calabi--Yau
threefolds with trivial canonical bundles and torsion-free second
integral homology. Fix a marking of the common complement of the node
neighborhoods and a relatively ample line bundle, restricting to $H_0$
and $H_S$. Write $C_i$ for the exceptional curves and set
\[
 W=\Bl_{\coprod_i C_i}Y,\quad \phi:W\to Y,\quad q=p\phi,
 \quad E=\sum_i E_i,\quad
 j=\PD_W^{-1}\phi^*\PD_Y.
\]
Here $j$ is an integral homology map, not the strict transform of a
chosen curve. Effective means represented by an effective integral
one-cycle, not merely belonging to a closed real cone.

\begin{proposition}\label{prop:finite}
The marking gives an integral surjection
\[
 \pi:H_2(Y;\ZZ)\longrightarrow H_2(S;\ZZ),\qquad
 \ker\pi=\langle[C_i]\rangle.
\]
The set
\begin{equation}\label{eq:F}
 F=\{\gamma\ne0:j\gamma\text{ is effective on }W\}
\end{equation}
is contained in $\Gamma_Y$, is additive after adjoining zero, contains
no class in $\ker\pi$, and has finitely many elements of bounded
$H_S\cdot\pi\gamma$ degree. The monoid
$\Gamma_T=\langle\Gamma_S\cup\pi(F)\rangle$ has positive, locally
finite degree $H_S\cdot(-)$.
\end{proposition}
\begin{proof}
Let $U$ be the common compact complement and
$L=\coprod_i(S^2\times S^3)$ its boundary. The local resolution pieces
retract to $S^2$, with boundary inclusion inducing an isomorphism on
$H_2$; the smoothing pieces retract to $S^3$. If
$a:H_2(L)\to H_2(U)$ is inclusion, integral Mayer--Vietoris gives
\[
 H_2(Y)=\bigl(H_2(U)\oplus\ZZ^r\bigr)/
       \{(a(v),-v):v\in\ZZ^r\}\simeq H_2(U),
 \qquad H_2(S)=H_2(U)/a(\ZZ^r).
\]
This proves the asserted quotient, including its integral meaning.

Locally a node is the affine cone over $\PP^1\times\PP^1$ with
polarization $\mathcal O(1,1)$. Blowing up its vertex gives
$\Tot(\mathcal O(-1,-1))$, as does blowing up the zero section of
either small resolution. These identifications glue, so
$W=\Bl_{\operatorname{Sing}X_0}X_0$. In particular $-E$ is $q$-ample,
and for some fixed integer $N>0$ the divisor
\[
 A=Nq^*H_0-E
\]
is ample. The projection formula gives $\phi_*j\gamma=\gamma$.
Restriction of $\phi^*\PD_Y\gamma$ to $E_i$ factors through
$H^4(C_i)=0$, hence $E_i\cdot j\gamma=0$. The marked polarization
agrees on $U$, and therefore
\begin{equation}\label{eq:ample-bound}
 A\cdot j\gamma=N H_S\cdot\pi\gamma.
\end{equation}
For $\gamma\in F$ the left side is positive. Pushforward proves
effectivity on $Y$, and sums of effective lifts prove additivity.
Bounded projected degree bounds ample degree on $W$. The homology
boundedness argument of Section~\ref{sec:universal}, and injectivity
of $j$, prove finiteness. Finally, the generators of $\Gamma_T$ of
degree at most $D$ form a finite set, and a sum of degree at most $D$
uses at most $D$ nonzero generators. This proves the final assertion.
\end{proof}

The integral blow-up decomposition also gives
\begin{equation}\label{eq:zero-contact}
 \operatorname{im}j=\bigcap_i\ker(E_i\cdot-).
\end{equation}
Indeed $H_2(W)=jH_2(Y)\oplus\bigoplus_i\ZZ a_i$, where $a_i$ is a
fiber of $E_i\to C_i$ and $E_i\cdot a_j=-\delta_{ij}$.

\begin{theorem}\label{thm:specialization}
Under the hypotheses of this section, put
\[
 \cE_Y=\prod_i f_0(Q^{[C_i]}),\qquad A_Y=\Z_Y/\cE_Y.
\]
Then $\Supp(A_Y-1)\subset F$. Coefficientwise summation along $\pi$
is finite in bounded target degree, and
\begin{equation}\label{eq:specialization}
 \Z_S=\Sp_\pi A_Y
   \quad\text{in }H\llbracket\Gamma_T\rrbracket.
\end{equation}
Here $\Z_S$ is extended by zero outside $\Gamma_S$, and
$[T^\beta]\Sp_\pi A_Y=\sum_{\pi\gamma=\beta}(A_Y)_\gamma$.
\end{theorem}

\begin{proof}
We first remove the entire exceptional contribution. The divisor
$p^*H_0$ is nef and vanishes precisely on contracted curves, so
$\Gamma_Y\cap\ker\pi$ is a face of $\Gamma_Y$, consisting of sums
$\sum m_i[C_i]$ with $m_i\ge0$. Restriction to this face commutes
with products, inverses and logarithms. In the resolved local node
$V=\Tot_{\PP^1}(\mathcal O(-1)^{\oplus2})$, the only compact
degree-$m$ cycle is $m[\PP^1]$: a compact curve maps to a point in the
affine node and must lie in the exceptional fiber. In Pardon's
Definition 4.7, one can take $N=0$ and no cut-down sections, giving
$(V,1_{m[\PP^1]})=x_{0,m}$. Locality and the product operation for
the disjoint node neighborhoods give
\[
 (Y,1_{\sum m_i C_i})=\prod_i x_{0,m_i}.
\]
At fixed class there are finitely many multiplicity vectors, bounded
by an ample degree on $Y$. Summing over them shows that the face
restriction of $\Z_Y$ is exactly $\cE_Y$, even when the exceptional
classes satisfy relations. Thus $A_Y$ restricts to $1$ on this face.

Next deform $Y$ to the normal cone of $\coprod C_i$. Its central
fiber is $W\cup_E\coprod_i\PP_{C_i}(N_{C_i/Y}\oplus\mathcal O)$.
For total relative class $\delta$ on $W$, contact order $\mu$ and
number of relative markings $\rho$, the virtual dimension is
\begin{equation}\label{eq:vdim}
 c_1(W)\cdot\delta+\rho-\mu=\rho-2\mu,
\end{equation}
since $K_W=E$. This holds also for disconnected $W$ parts by summing
over vertices. If contacts are nonempty, $\mu\ge\rho>0$, and the
virtual class vanishes. A surviving connected noncontracted term
lies entirely on $W$; a term on a projective bundle pushes to a
multiple of $[C_i]$. Equation~\eqref{eq:zero-contact} therefore gives
\begin{equation}\label{eq:relative}
 N^Y_{h,\gamma}=N^{(W,E)}_{h,j\gamma\mid\varnothing}
      \qquad(\pi\gamma\ne0).
\end{equation}
This is the unmarked resolution comparison used by Lee--Lin--Wang
\cite[Section 2.1]{LLW}. In particular all genera vanish when
$\pi\gamma\ne0$ and $\gamma\notin F$. For $L_Y=\log A_Y$, connected evaluation at a
noncontracted class equals $N^Y_{h,\gamma}$; at a contracted nonzero
class $L_Y$ itself is zero. The support consequence of
Lemma~\ref{lem:formal} gives $\Supp L_Y\subset F$, and additivity
of $F\cup\{0\}$ gives $\Supp(A_Y-1)\subset F$.

Proposition~\ref{prop:finite} now licenses all sums in $\Sp_\pi A_Y$,
including products. Coproduct commutes with the finite sums. For
division, only $\gamma=d\eta$ contributes, and torsion-freeness of
the target gives
\[
 \rho_d\sum_{\pi\gamma=\beta}(A_Y)_\gamma
 =\sum_{\pi\eta=\beta/d}(A_Y)_\eta.
\]
Thus the specialization is compatible. The zero extension of $\Z_S$
is compatible as well: sums and positive multiples of effective
classes remain effective.

It remains to compare connected evaluations also at classes of
$\Gamma_T$ which need not be effective on $S$. Semistable reduction
of the smoothing gives a smooth total space $\mathscr V$ with central
fiber $W\cup_E\coprod_i Q_i$, where $Q_i$ is a smooth quadric
threefold and $E_i$ is its hyperplane section
\cite[Section 3.1]{LW}. If $i:S\to\mathscr V$ and
$i_0:W\to\mathscr V$ are inclusions, the common-complement marking
gives
\begin{equation}\label{eq:total-class}
 i_{0*}j\gamma=i_*\pi\gamma.
\end{equation}
Indeed a representative in $U$ has zero exceptional intersections on
$W$, hence represents $j\gamma$ by \eqref{eq:zero-contact}; its
transport through the smooth complement supplies the required
cylinder. Node monodromy is trivial on $H_2(S)$, and the invariant
cycle theorem gives rational injectivity of $i_*$, as in
\cite[Lemmas 2.18 and 3.4]{LW}. Its integral kernel is torsion and
therefore zero here.

For every $\beta\in\Gamma_T$, the class $i_*\beta$ is effective in
$\mathscr V$: this holds for generators from $S$, and for generators
from $F$ by \eqref{eq:total-class}, then for their sums. Apply the
degeneration formula graded by this total-space class
\cite[Theorem 2.16]{LW}, retaining connected dual graphs.
This grading restricts the virtual-cycle formula to open-and-closed
components of the map moduli space, so its smoothing side may be
empty. It is zero when $\beta$ is ineffective on $S$; injectivity of
$i_*$ otherwise identifies its sole class with $\beta$.
Nonempty contacts again vanish by \eqref{eq:vdim}. A nonzero
empty-contact curve on $Q_i$ is impossible because $E_i$ is ample.
Surviving classes on $W$ are exactly $j\gamma$ with
$\gamma\in F$ and $\pi\gamma=\beta$. Consequently, for all
$0\ne\beta\in\Gamma_T$ and every $h$,
\[
 N^S_{h,\beta}=\sum_{\substack{\gamma\in F\\\pi\gamma=\beta}}
 N^{(W,E)}_{h,j\gamma\mid\varnothing}
 =\sum_{\substack{\gamma\in F\\\pi\gamma=\beta}}N^Y_{h,\gamma}.
\]
These are precisely the connected evaluations of the two compatible
series in \eqref{eq:specialization}. Lemma~\ref{lem:formal} proves
their equality. In particular no preservation of effectivity under
smoothing was assumed.
\end{proof}

\begin{corollary}\label{cor:integral}
The rational product exponents satisfy
\[
 e_{\beta,g}(S)=\sum_{\substack{\gamma\in F\\\pi\gamma=\beta}}
 e^{\mathrm{red}}_{\gamma,g}(Y).
\]
In particular an integral product expansion on $Y$ implies one on $S$.
\end{corollary}
\begin{proof}
Take logarithms in \eqref{eq:specialization} and compare the
$p_{g,1}$ coefficients in Lemma~\ref{lem:formal}. Normalization
changes only exceptional exponents, so the reduced and original
exponents agree at every noncontracted class. All sums are finite.
\end{proof}

There is a useful chamber independence in this formulation. If $Y_-$
and $Y_+$ are two projective small resolutions of the same nodal
threefold, their blow-ups are the same $W$. The maps $j_\pm$ have
common image \eqref{eq:zero-contact}, so
$\Phi=\phi_{+*}j_-$ is an integral isomorphism satisfying
$j_+\Phi=j_-$ and $\Phi(F_-)=F_+$. Their normalized series agree
under this identification: \eqref{eq:relative} identifies all
connected evaluations in the common positive completion, and
reconstruction applies. A flopped exceptional class changes sign;
this statement does not identify the unnormalized effective expansions.
